% Bài làm buổi 18/11
\begin{problem}[8DS-HSBN-005]
Đồ thị song song với đường thẳng $y = -3x + 5$ nên $a = -3$ và $b \neq 5$.

Đồ thị cắt trục hoành tại điểm có hoành độ bằng $2$ nên đi qua điểm $(2; 0)$: $0 = a \cdot 2 + b$.

Thay $a = -3$: $0 = -6 + b$, suy ra $b = 6$.

Đối chiếu: $b = 6 \neq 5$ (thoả mãn), $a = -3 \neq 0$ nên hàm số tìm được là hàm số bậc nhất.

Đáp số: $y = -3x + 6$
\end{problem}

\begin{problem}[8DS-PTNT-004]
a) Tấm bìa có diện tích $a^2$ (cm$^2$), mỗi hình vuông nhỏ có diện tích $b^2$ (cm$^2$), nên $S = a^2 - 4b^2$ (cm$^2$).

$S = a^2 - (2b)^2 = (a - 2b)(a + 2b)$.

b) Với $a = 46$, $b = 2$ (thoả mãn $0 < 2b < a$): $S = (46 - 4)(46 + 4) = 42 \cdot 50 = 2\,100$.

Đáp số: a) $S = (a - 2b)(a + 2b)$; b) $S = 2\,100$
\end{problem}
